\section{Discussion}\label{sec:discussion}

\subsection{Limitations}
\paragraph{Hot-row contention} On \texttt{hot-row-update} at 64
threads RedlineDB lands at 0.21$\times$ SQLite's throughput
(Table~\ref{tab:headline}). The cause is structural: when every
contender writes the same row, the row-lock manager queues all but
one writer at a time, and there is no shape of MVCC that helps. SQLite
batches concurrent updates through its WAL writer, amortizing the
fsync; we currently do not coalesce row-lock holders that target
identical rows. A future commit-batching pass that detects
hot-row contenders and applies their effects in a single
group-commit cycle would close most of the gap.

\paragraph{Secondary-index range scans} On
\texttt{secondary-index-range} at 64 threads, RedlineDB is roughly
80$\times$ slower than SQLite. The proximate cause is that our
index cursor walks one leaf at a time without prefetch. SQLite's
range scan reads ahead while the consumer is processing. Adding a
prefetch ring (one outstanding page-read in flight per cursor) is
known engineering work; it does not require a kernel-shape change.

\paragraph{Single-thread per-tx overhead} The fixed overhead of a
single-threaded transaction is higher in RedlineDB than in SQLite
because we route through the row-lock manager and CSN allocator
even when no contention is present. The 1-thread point-read row in
Table~\ref{tab:headline} is 3.22$\times$ \emph{faster} than SQLite,
so this overhead is not user-visible at the threading regimes we
target; but for a single-threaded read-heavy use case where SQLite
is already saturating the host, the overhead is real.

\paragraph{No encryption-at-rest} RedlineDB does not currently
encrypt page data on disk. The kernel structure is friendly to a
page-level encryption layer between the page cache and the page
file---encryption is a known follow-on, not a re-architecture---but
shipping it is future work.

\subsection{Future Work}
The natural next step on the concurrency axis is serializable
snapshot isolation~\cite{cahill2008ssi}: SI plus a runtime
dangerous-structure detector. The certification protocol already
has the determinism it needs to validate an SSI implementation
against the workload matrix. On the executor axis, a vectorized
batch-tuple path along the lines of
HyPer~\cite{kemper2011hyper,neumann2015fast} would close the
analytical workload gap without intruding on the OLTP path. On the
ecosystem axis, a libSQL~\cite{libsql} migration tool---a small
adapter that translates a libSQL connection string to a RedlineDB
connection string---would let existing libSQL clients try the engine
with no code changes.
